\section{Verification purpose}
This case verifies dynamic herbaceous curing in GR2 and GS3.  Live-herb moisture is the controlled input; ELMFIRE must transfer the corresponding fraction of herbaceous load to the dead 1-h class and then evaluate the full Rothermel model.  It can expose reversed curing, incorrect 30/120\% bounds, failure to conserve load, wrong class SAV weighting, or confusion between curing and live moisture damping.  The verification question is whether all required curing states reproduce independently calculated intensity and ROS within 1\% and \RosTolerance{}, respectively.

\section{Mathematical formulation and numerical implementation}
For live-herb moisture $M_{LH}$ expressed as a fraction, the cured fraction is
\begin{equation}T=\max[0,\min(1,1.33-1.11M_{LH})].\end{equation}
ELMFIRE implements the equivalent integer-percent interpolation between 30 and 120\%.  If $W_{LH,0}$ is original live-herb load, $W_{1h}'=W_{1h}+TW_{LH,0}$ and $W_{LH}'=(1-T)W_{LH,0}$.  The transferred material adopts dead 1-h moisture while remaining live material uses live damping and heat of preignition.  The formulation follows Andrews (2018), RMRS-GTR-371.

\section{Derivation of expected behavior}
At 30\% moisture all live-herb load is cured; at 75\% half is cured; at 120\% and above none is transferred.  Above 120\%, curing is constant but live moisture damping continues to change, so ROS need not form a plateau.  The case-local oracle conserves total load, recomputes effective SAV and every dependent Rothermel coefficient, applies moisture damping, and applies ELMFIRE's selected wind limit at the target \SI{6}{mph} midflame wind.

\section{Verification metrics and acceptance criteria}
Maximum relative errors are 0.02 for direct ROS and 0.01 for reaction intensity.  Variants predicted at or above \SI{0.10}{m.min^{-1}} also require TOA-derived ROS error at most 0.05 and minimum $R^2=0.995$; the most moisture-limited variants use direct fields because they cannot cross a multi-cell corridor within the practical duration cap.  Direct effective loads are not exposed by the normal executable, so their conservation is an oracle invariant and the implementation is observed through intensity and ROS.  All 20 variants are required; no qualitative trend alone can pass the case.

\section{Simulation configuration and predicted outcome}

The finite-strip PHI raster is a signed-distance field clipped to ELMFIRE's supported $[-1,1]$ interval, preserving both an ignited interior and a resolved front gradient. For every positive oracle ROS, \texttt{SIMULATION\_DT} and \texttt{SIMULATION\_DTMAX} are selected so the expected front advances 5\% of one cell per step (CFL 0.05, below the configured 0.2 limit). This keeps the PHI change above ELMFIRE's fixed stall scale even if the implementation produces only a small fraction of the reference speed. A zero-ROS variant retains a 5 s step and is evaluated from direct fields without requiring TOA propagation. The analytically selected stop time is rounded upward to a whole second because the current ELMFIRE time loop stores the duration as an integer; this preserves the intended duration and guarantees the terminal output dump. The evaluator always requires uniquely named rasters tied to the one final dump record. A TOA-required variant must reach the \texttt{SIMULATION\_TSTOP} in its generated namelist; a direct-only variant may terminate early because lack of propagation is part of the state being tested.
Live-herb moisture is 30, 45, 60, 75, 90, 100, 110, 120, 135, or 200\%.  Dead classes are fixed at 5\%, live woody moisture at 90\%, slope at zero, and target midflame wind at \SI{6}{mph}; the 20-ft input is back-calculated from each fuel's WAF.  Each homogeneous domain uses a finite ignition strip with a planar east face inside the analysis corridor and with lateral ends clear of the grid boundary.  Crown fire, spotting, buildings, and acceleration are excluded.  Strong nonlinear response is expected in the curing interval.

\section{Actual simulation results}
The generated record contains \CompletedVariantCount{} of \RequiredVariantCount{} variants and reports \textbf{\VerificationStatus}.
\ifnum\CompletedVariantCount>0
\IfFileExists{../figures/sweep_response.pdf}{\begin{figure}[H]\centering\includegraphics[width=0.95\textwidth]{sweep_response.pdf}\caption{Hand calculations (open diamonds, solid lines) and ELMFIRE predictions (filled circles, dashed lines) for the dynamic-curing response. Each fuel-model pair uses the same color.}\end{figure}}{}
\fi

\section{Calculated metrics and verification decision}
\begin{center}\begin{tabular}{@{}llll@{}}\toprule Metric & Limit & Calculated & Status\\\midrule\MetricRows\bottomrule\end{tabular}\end{center}
The authoritative decision is \textbf{\VerificationStatus}.

\section{Technical interpretation}
A reversed curve indicates a reversed transfer fraction.  Discontinuities away from integer moisture selection suggest class indexing or interpolation defects.  Correct behavior above 120\% but disagreement below it isolates curing rather than live damping.  Correct intensity and wrong ROS implicate heat sink or wind coupling.

\section{Scope and limitations}
This executable-level test observes curing indirectly; a direct internal-load assertion would require a compiled component diagnostic.  It covers two dynamic fuel models, constant weather, and no temporal curing.  CASE38 separately isolates live moisture without load transfer.
